\documentclass[11pt]{article}
\usepackage[T1]{fontenc}\usepackage[utf8]{inputenc}
\usepackage[spanish,es-noindentfirst]{babel}\usepackage{lmodern}\usepackage{microtype}
\usepackage[letterpaper,margin=1in]{geometry}\usepackage{parskip}
\usepackage{enumitem}\usepackage{xcolor}\usepackage{mdframed}
\usepackage[hidelinks]{hyperref}\usepackage{xurl}
\definecolor{qbg}{HTML}{F4F1EA}\definecolor{qrule}{HTML}{7A6A53}
\definecolor{ctx}{HTML}{6B6B6B}
\newcommand{\score}[2]{{\small\textbf{Chroma score:} #1 \quad|\quad \textbf{chunk} \##2}}
\setlength{\parskip}{0.5\baselineskip}
\title{\vspace{-2em}\bfseries Decálogo de ética profesional docente\\[2pt]\large Cuatro principios por artículo, anclados en Chroma}
\author{Jesús G.}\date{11 de julio de 2026}
\begin{document}\maketitle\thispagestyle{empty}
\begin{quote}\small\textit{Nota metodológica.} Cada punto se deriva del artículo indicado y se contrasta con su texto completo mediante un índice vectorial local (Chroma, similitud coseno). Se reporta el \emph{score}, el \emph{chunk} que respalda el punto (cita textual) y un fragmento del chunk anterior y del posterior como contexto. Nota: los dos archivos \texttt{4406374} entregados son idénticos (mismo md5), por lo que se tratan como un solo artículo.\end{quote}
\vspace{0.4em}
\section*{Art\'iculo 1 --- \'Etica profesional basada en principios (Hirsch Adler, 2013)}
{\small\textit{Fuente: }Hirsch Adler, A. (2013). La \'etica profesional basada en principios y su relaci\'on con la docencia. \textit{Edetania, 43}, 97--111.}\par\vspace{0.4em}
\noindent\textbf{Punto 1. Beneficencia.} Hacer el bien propio de la profesi\'on: en la docencia, el bien intr\'inseco es que los alumnos aprendan, entendiendo la ense\~nanza como ayudar a aprender.\par
\score{0.702}{34}\par\smallskip
\noindent{\footnotesize\color{ctx}\textit{\ldots{}chunk anterior (\#33): }bienes intrínsecos EDETANIA 43 [Julio 2013], 97-111, ISSN: 0214-8560 Ana Hirsch Adler 106 son los prioritarios y que se tergiversan las actividades profesionales cuando los esfuerzos están dirigidos únicamente al logro de beneficios personales. Como\ldots{}}\par\smallskip
\begin{mdframed}[backgroundcolor=qbg,linecolor=qrule,linewidth=1pt,leftline=true,topline=false,bottomline=false,rightline=false,innerleftmargin=8pt,innertopmargin=6pt,innerbottommargin=6pt]
{\small\textbf{Cita textual (chunk \#34):} «ya se había mencionado en la introducción, el Dr. Hortal retoma los principios generales de la ética profesional para reflexionar en torno a la docencia. Con respecto al principio de beneficencia, considera que: El bien intrínseco de la práctica de la docencia es que los alumnos aprendan. El ejercicio éticamente responsable de la función docente lleva consigo al menos estos deberes y responsabilidades: ante todo, enseñar, entendiendo la enseñanza como ayudar a aprender. Enseñar presupone saber, haber aprendido lo que se enseña y estar al día en la materia que se enseña (Hortal, 2000: 61). Resulta importante completar esa idea del Dr. Hortal con lo expresado al respecto por Rafaela»}
\end{mdframed}
\noindent{\footnotesize\color{ctx}\textit{chunk posterior (\#35): }García-López, Gonzalo Jover y Juan Escámez (2010: 17-19), acerca de que ``la meta social de la docencia consiste en la transmisión de la cultura y la formación de personas críticas''. Al retomar a Fernando Savater, indican que se trata de\ldots{}}\par
\vspace{0.9em}
\noindent\textbf{Punto 2. No maleficencia.} Por sobre todo, no hacer da\~no: evitar el da\~no f\'isico, emocional o a la autoestima del estudiante, y no dejarlo en peor situaci\'on que antes.\par
\score{0.552}{37}\par\smallskip
\noindent{\footnotesize\color{ctx}\textit{\ldots{}chunk anterior (\#36): }que desconoce en el campo que se está tratando (García-López et. al., 2010: 19). Principio de no maleficencia El principio de no maleficencia enfatiza la obligación3 de no infligir daño a otros: ``por sobre todo\ldots{}}\par\smallskip
\begin{mdframed}[backgroundcolor=qbg,linecolor=qrule,linewidth=1pt,leftline=true,topline=false,bottomline=false,rightline=false,innerleftmargin=8pt,innertopmargin=6pt,innerbottommargin=6pt]
{\small\textbf{Cita textual (chunk \#37):} «no hacer daño''. Consiste en actuar de manera que no se ponga en ries3 El juramento hipocrático expresa claramente la obligación de no maleficencia y la de beneficencia: ``Usaré el tratamiento para ayudar a los enfermos de acuerdo a mi habilidad y mi juicio y nunca los usaré para perjudicarlos'' (Beuchamp y Childress, 2001: 113). EDETANIA 43 [Julio 2013], 97-111, ISSN: 0214-8560 La ética profesional basada en principios y su relación con la docencia 107 go o se lastime a las personas. Aunque este principio se utiliza en mayor medida en las ciencias médicas y en la ética de la investigación científica, consideramos que es también relevante en el marco»}
\end{mdframed}
\noindent{\footnotesize\color{ctx}\textit{chunk posterior (\#38): }de todas las profesiones, incluida la profesión docente. La protección contra el daño tiene al menos dos vertientes (David y Sutton, 2011): a) Del daño físico. Aunque esta posibilidad es más común en la investigación médica y biológica, puede producirse\ldots{}}\par
\vspace{0.9em}
\noindent\textbf{Punto 3. Autonom\'ia.} Reconocer la capacidad de darse las propias normas: el docente requiere independencia para actuar \'eticamente y debe respetar la autonom\'ia del estudiante.\par
\score{0.670}{41}\par\smallskip
\noindent{\footnotesize\color{ctx}\textit{\ldots{}chunk anterior (\#40): }libre de interferencias de control por parte de otros y de contar con un entendimiento adecuado para tomar decisiones significativas (capacidad para la acción intencionada). 4 Beauchamp y Chidress (2001) citan varios casos reales que\ldots{}}\par\smallskip
\begin{mdframed}[backgroundcolor=qbg,linecolor=qrule,linewidth=1pt,leftline=true,topline=false,bottomline=false,rightline=false,innerleftmargin=8pt,innertopmargin=6pt,innerbottommargin=6pt]
{\small\textbf{Cita textual (chunk \#41):} «se produjeron en Estados Unidos al respecto. EDETANIA 43 [Julio 2013], 97-111, ISSN: 0214-8560 Ana Hirsch Adler 108 Autonomía del beneficiario En el segundo caso (que es la propuesta de Beauchamp y Childress, 2001; Hortal, 2002, y Bermejo, 2002), el principio de autonomía busca corregir la falta de simetría entre quien ofrece el servicio y el beneficiario de la actividad. El profesional por su preparación, acreditación y dedicación tiene un ascendente sobre sus clientes y usuarios. La desigualdad entre ambas partes puede producir abusos (entre ellos el paternalismo). Para evitarlos, es necesario que esté siempre en funcionamiento el principio de autonomía. Consiste en considerar que el receptor de los servicios»}
\end{mdframed}
\noindent{\footnotesize\color{ctx}\textit{chunk posterior (\#42): }(individual y colectivo) no es un ente pasivo, sino un sujeto protagonista. De ahí, se deriva la obligación de garantizar a todos los individuos involucrados el derecho de ser informados, de que se respeten sus derechos y de consentir antes\ldots{}}\par
\vspace{0.9em}
\noindent\textbf{Punto 4. Justicia.} Distribuir de manera equitativa los beneficios y las obligaciones, con procedimientos razonables y no explotadores.\par
\score{0.646}{13}\par\smallskip
\noindent{\footnotesize\color{ctx}\textit{\ldots{}chunk anterior (\#12): }informado. 1. ``Es absolutamente esencial el consentimiento voluntario del sujeto humano''. Esto significa que la persona implicada debe tener capacidad legal para dar consentimiento; su situación debe ser tal que pueda ser capaz de ejercer\ldots{}}\par\smallskip
\begin{mdframed}[backgroundcolor=qbg,linecolor=qrule,linewidth=1pt,leftline=true,topline=false,bottomline=false,rightline=false,innerleftmargin=8pt,innertopmargin=6pt,innerbottommargin=6pt]
{\small\textbf{Cita textual (chunk \#13):} «una elección libre, sin intervención de cualquier elemento de fuerza, fraude, engaño, coacción u otra forma de constreñimiento o coerción; debe tener suficiente conocimiento y comprensión de los elementos implicados que le capaciten para hacer una decisión razonable e ilustrada. Este último elemento requiere que antes de que el sujeto de experimentación acepte una decisión afirmativa, debe conocer la naturaleza, duración y fines del experimento, el método y los medios con los que será realizado, todos los inconvenientes y riesgos que pueden ser esperados razonablemente y los efectos sobre su salud y persona que pueden posiblemente originarse de su participación en el experimento. El deber y la responsabilidad para asegurarse»}
\end{mdframed}
\noindent{\footnotesize\color{ctx}\textit{chunk posterior (\#14): }de la calidad del consentimiento residen en cada individuo que inicie, dirija o esté implicado en el experimento. Es un deber y responsabilidad personales que no pueden ser delegados impunemente. EDETANIA 43 [Julio 2013], 97-111, ISSN: 0214-8560 La ética profesional\ldots{}}\par
\vspace{0.9em}
\section*{Art\'iculo 2 --- \'Etica profesional docente (Rojas Artavia, 2011)}
{\small\textit{Fuente: }Rojas Artavia, C. E. (2011). \'Etica profesional docente: un compromiso pedag\'ogico human\'istico. \textit{Revista Humanidades, 1}, 1--22.}\par\vspace{0.4em}
\noindent\textbf{Punto 1. M\'as que la norma.} La \'etica profesional docente no se reduce a la deontolog\'ia ni al cumplimiento de la norma escrita: es un compromiso vivencial que va m\'as all\'a de ella.\par
\score{0.742}{26}\par\smallskip
\noindent{\footnotesize\color{ctx}\textit{\ldots{}chunk anterior (\#25): }es primero profesional que sujeto moral. Por ello, es que la moral vivida forma parte habitual e irrecusable de nuestras actuaciones, en cualquier esfera social en que nos encontremos. Se es profesional al ostentar -\ldots{}}\par\smallskip
\begin{mdframed}[backgroundcolor=qbg,linecolor=qrule,linewidth=1pt,leftline=true,topline=false,bottomline=false,rightline=false,innerleftmargin=8pt,innertopmargin=6pt,innerbottommargin=6pt]
{\small\textbf{Cita textual (chunk \#26):} «además del conocimiento especializado para el ejercicio de la labor- principios éticos que nos hacen ser personas libres y responsables de nuestras acciones e imputables moralmente por ellas desde una normativa codificada que compete al gremio al cual se pertenece. En congruencia con lo anterior, se puede decir que la profesión como práctica social, guarda un fundamento eminentemente iluminado por el bien común, tal y como apunta Muñoz: ``El código de moral profesional --inmerso en el código moral más amplio de la sociedadcrea como expectativa más importante la de un ejercicio profesional óptimo y responsable en beneficio del bien común (...) el profesional no sólo desempeña una función social que»}
\end{mdframed}
\noindent{\footnotesize\color{ctx}\textit{chunk posterior (\#27): }contribuye al bien común (...) también, y más específicamente, cada profesional forma parte de la intelligentsia de la comunidad del país. La totalidad de los profesionales son la intelligentsia y son los usufructuarios, promotores y depositarios del saber logrado por\ldots{}}\par
\vspace{0.9em}
\noindent\textbf{Punto 2. Educar es humanizar.} La finalidad de la educaci\'on es la formaci\'on integral ---cognitiva, afectiva y moral--- de la persona; educar es humanizar, no acumular datos ni s\'olo capacitar en una t\'ecnica.\par
\score{0.698}{1}\par\smallskip
\noindent{\footnotesize\color{ctx}\textit{\ldots{}chunk anterior (\#0): }un compromiso vivencial que va más allá de la norma escrita y debe hacerse efectivo teórica y prácticamente. En el ámbito de la educación, para cumplir con ese compromiso el y la docente han de\ldots{}}\par\smallskip
\begin{mdframed}[backgroundcolor=qbg,linecolor=qrule,linewidth=1pt,leftline=true,topline=false,bottomline=false,rightline=false,innerleftmargin=8pt,innertopmargin=6pt,innerbottommargin=6pt]
{\small\textbf{Cita textual (chunk \#1):} «ser conscientes de sus tenencias prácticas, intelectivas y morales, así como del deber de desarrollarlas constantemente para ponerlas a disposición de sus estudiantes y ayudarles a crecer cognitiva, afectiva y moralmente de manera integral. Palabras clave: Ética profesional - docencia - tenencia práctica, intelectiva y moral -- compromiso y virtud. Abstract The professional ethics is not simply a deontology or a group of norms to govern the behavior of those who exercise a professional work, it is a commitment that goes beyond the written norm and it should become effective theoretical and practically. In the environment of education, to fulfill that commitment the educators must be aware of their practical,»}
\end{mdframed}
\noindent{\footnotesize\color{ctx}\textit{chunk posterior (\#2): }intellectual and moral possessions, as well as their duty of constantly developing them to put them to their student's disposition and help them grow cognitive, affective and morally in an integral way. Key words: Professional Ethics -- education- cognitive, intellectual\ldots{}}\par
\vspace{0.9em}
\noindent\textbf{Punto 3. Naturaleza social e hist\'orica.} El ser humano es social, hist\'orico y din\'amico: no sobrevive aislado, necesita de los otros y de la cultura, que la educaci\'on transmite de generaci\'on en generaci\'on.\par
\score{0.747}{18}\par\smallskip
\noindent{\footnotesize\color{ctx}\textit{\ldots{}chunk anterior (\#17): }tirantez constante entre el ser (su concretitud) y el deber ser (ideal por alcanzar). En términos filosóficos, el ser humano es el resultado constante y cambiante de ese proceso dialéctico entre lo que se es\ldots{}}\par\smallskip
\begin{mdframed}[backgroundcolor=qbg,linecolor=qrule,linewidth=1pt,leftline=true,topline=false,bottomline=false,rightline=false,innerleftmargin=8pt,innertopmargin=6pt,innerbottommargin=6pt]
{\small\textbf{Cita textual (chunk \#18):} «y lo que se pretende ser. Además, la naturaleza humana es histórica: Se sobrevive a partir de la reproducción de aquello que la cultura -entendida ésta como todo aquello hecho por el ser humano - considera adecuado para la prosecución de la vida. El cambio constante en sus condiciones de vida le conduce a mantener un acumulado histórico de experiencias que, como necesidad -- capacidad de adaptación, tanto al medio natural como social le lleva a hacer operativo ese acervo socio cultural en pos, no sólo de su sobrevivencia sino de una mayor calidad de vida. En este sentido, la cultura logra sus propósitos comisionándole a la educación --formal e»}
\end{mdframed}
\noindent{\footnotesize\color{ctx}\textit{chunk posterior (\#19): }informal- el papel de trasmitir de generación en generación esta experiencia histórica acumulada, que muta según se gesten nuevas necesidades. El ser humano conserva su pasado, lo torna operativo en el presente para proyectar su vida a futuro. De esta\ldots{}}\par
\vspace{0.9em}
\noindent\textbf{Punto 4. Cuidado del otro y ethos.} La acci\'on educativa se funda en la responsabilidad y el cuidado del otro; el/la docente se vuelve modelo por la integridad que manifiesta en su hacer.\par
\score{0.664}{78}\par\smallskip
\noindent{\footnotesize\color{ctx}\textit{\ldots{}chunk anterior (\#77): }un código de ética profesional --que lo debe haber- , sino porque se cree profundamente en lo que se hace. No se debe actuar moralmente sólo porque hay mecanismos externos que nos compulsan, sino por\ldots{}}\par\smallskip
\begin{mdframed}[backgroundcolor=qbg,linecolor=qrule,linewidth=1pt,leftline=true,topline=false,bottomline=false,rightline=false,innerleftmargin=8pt,innertopmargin=6pt,innerbottommargin=6pt]
{\small\textbf{Cita textual (chunk \#78):} «respeto y deber para con uno mismo y para con el otro. El ser humano no es un simple individuo al que se debe proveer de un conjunto de capacidades técnico procedimentales para que saque provecho de sus habilidades; se trata de que unido al afinamiento de esas capacidades se le fomente el compromiso con su entorno y con la responsabilidad histórica que le compete. Si se ha establecido que por naturaleza el ser humano es social, histórico y dinámico, por esa misma condición ineludible, necesita del otro, no sólo para sobrevivir, sino para convivir. El compromiso de los y las docentes no es para con una idea abstracta o»}
\end{mdframed}
\noindent{\footnotesize\color{ctx}\textit{chunk posterior (\#79): }``técnica'' del ser humano, sino para con el otro, para con ese ser concreto que existe y me interpela como rostro --como diría Levinas- y al que hay que escuchar y cuidar mediante nuestro saber y hacer educativo. No se\ldots{}}\par
\vspace{0.9em}
\end{document}